\section{Desarrollo del Proyecto}
El trabajo se dividió en cuatro lotes, uno por integrante: la operativa de empresas, el backoffice y la integración con balanzas; la aplicación móvil, la identidad y el despliegue; la ingesta de eventos, los nodos periféricos, la performance y la observabilidad; y la fiscalización automática. Las decisiones de arquitectura que afectaban a más de un lote se discutieron en conjunto y se registraron en el documento de arquitectura.

El código se gestionó en un repositorio Git. Cada cambio se desarrolló en una rama propia y se integró a la rama principal mediante un pull request; entre el 6 de octubre y el 2 de noviembre se integraron 81. Un pipeline de integración continua compila el componente central, ejecuta las pruebas automatizadas y valida el frontoffice y la aplicación móvil en cada cambio, y desde una rama de release se despliegan automáticamente el frontoffice y el componente central. El documento de arquitectura se mantuvo en LaTeX en el mismo repositorio, lo que permitió revisar sus cambios con el mismo proceso que el código.

El seguimiento se realizó mediante monitoreos semanales con el docente. Las entregas intermedias se corrieron una semana respecto de las fechas previstas en la letra. La principal desviación fue el despliegue: la plataforma Elastic Cloud prevista para el componente central no estuvo disponible, por lo que se buscó una alternativa.

\textcolor{red}{[Completar: porcentaje de tiempo dedicado a cada tarea y comparación con el cronograma inicial.]}
